%% TU Delft House Colors
\definecolor{tudelft-blue}{cmyk}{0.75,0.16,0.05,0}
\definecolor{tudelft-white}{cmyk}{0,0,0,0}
\definecolor{tudelft-black}{cmyk}{0,0,0,1}

% TU Delft Secondary colors
\definecolor{tudelft-darkblue}{cmyk}{1.0,0.8,0.08,0.7}
\definecolor{tudelft-lightcyan}{cmyk}{0.72,0, 0.24,0.0}
\definecolor{tudelft-lightblue}{cmyk}{0.98,0.4,0,0}
\definecolor{tudelft-purple}{cmyk}{0.65,1.0,0,0.12}
\definecolor{tudelft-pink}{cmyk}{0,0.7,0,0}
\definecolor{tudelft-red}{cmyk}{0.05,1,0.48,0.3}
\definecolor{tudelft-darkorange}{cmyk}{0,0.85,0.75,0}
\definecolor{tudelft-orange}{cmyk}{0,0.7,0.75,0}
\definecolor{tudelft-yellow}{cmyk}{0,0.31,0.98,0}
\definecolor{tudelft-lightgreen}{cmyk}{0.63,0,0.84,0}
\definecolor{tudelft-darkgreen}{cmyk}{1,0,0.68,0.04}

% Choose primary color
\colorlet{primary}{tudelft-darkblue}

% Choose secondary color
\colorlet{secondary}{tudelft-blue}

% Choose accent color
\colorlet{accent}{tudelft-red}

% _____ CALLOUT BOXES _____

% tcolorbox is required here rather than in main.tex because this file is
% \input by the class (cls:87), so it must be loaded by this point for the
% definitions below to see it.
\RequirePackage[most]{tcolorbox}

% The contributions/discoveries callout that closes a chapter. Usage:
%     \begin{contributions}                  % default title
%         \item \textbf{Lead-in.} Body text.
%     \end{contributions}
%     \begin{contributions}[Another title]   % title overridden
\newtcolorbox{contributionsbox}[1]{
    enhanced, breakable, sharp corners,
    colback=primary!4, colframe=primary, boxrule=0.9pt,
    title={#1}, fonttitle=\bfseries, coltitle=white,
    left=5mm, right=5mm, top=4mm, bottom=4mm,
}

\newenvironment{contributions}[1][Contributions and discoveries]
    {\begin{contributionsbox}{#1}%
     \begin{itemize}[leftmargin=5.5mm, itemsep=1.6ex, topsep=0.5ex, parsep=0pt,
                     label=\textcolor{primary}{\rule[0.35ex]{3.6pt}{3.6pt}}]}
    {\end{itemize}\end{contributionsbox}}

% The key-idea callout that closes a section. Deliberately lighter than the
% contributions box above, since several appear in one chapter: white body, a
% hairline frame and a tinted title bar rather than a solid one. Usage:
%     \begin{keyideas}                     % title reads "Key ideas"
%         \item \textbf{Lead-in.} Body text.
%     \end{keyideas}
%     \begin{keyideas}[Timescales]         % title reads "Timescales"
\newtcolorbox{keyideasbox}[1]{
    enhanced, breakable, arc=0.8mm,
    colback=white, colframe=primary!55, boxrule=0.4pt,
    colbacktitle=primary!8, coltitle=primary, fonttitle=\bfseries,
    title={#1},
    left=4mm, right=4mm, top=3mm, bottom=3mm,
    before skip=3.5mm, after skip=3.5mm,
}

\newenvironment{keyideas}[1][Key ideas]
    {\begin{keyideasbox}{#1}%
     \begin{itemize}[leftmargin=5mm, itemsep=1.2ex, topsep=0.3ex, parsep=0pt,
                     label=\textcolor{primary!70}{\rule[0.35ex]{3pt}{3pt}}]}
    {\end{itemize}\end{keyideasbox}}

% _____ PRIMITIVE STYLES _____

\tikzstyle{block} = [
    draw,
    minimum width=2.5cm,
    minimum height=1.5cm,
    rectangle, 
    rounded corners=4, 
    line width=1.2,
    text=black,
    align=center,
]

\tikzstyle{line} = [
    draw, 
    -latex'
]

\tikzstyle{line-primary} = [
    draw=primary, 
    -latex'
]

\tikzstyle{line-secondary} = [
    draw=secondary, 
    -latex'
]

\tikzstyle{line-accent} = [
    draw=accent, 
    -latex'
]

\tikzstyle{thickline} = [
    line,
    ultra thick
]

% ____ Colored primitives ____

\tikzstyle{block-primary} = [
    block,
    draw=primary
]

\tikzstyle{block-secondary} = [
    block,
    draw=secondary
]

\tikzstyle{block-accent} = [
    block,
    draw=accent
]

\tikzstyle{line-primary} = [
     line,
     draw=primary
]


\tikzstyle{line-secondary} = [
     line,
     draw=secondary
]


\tikzstyle{line-accent} = [
     line,
     draw=accent
]

% ____ Shaded primitives _____

\tikzstyle{shaded-pattern} = [
    pattern = north east lines
]

\tikzstyle{shaded-primary} = [
    block-primary,
    shaded-pattern,
    pattern color=primary!30,
]

\tikzstyle{shaded-secondary} = [
    block-secondary,
    shaded-pattern,
    pattern color=secondary!40,
]

\tikzstyle{shaded-accent} = [
    block-accent,
    shaded-pattern,
    pattern color=accent!40,
]

% Branch and junction points

\tikzstyle{branch} = [
    fill,
    shape=circle,
    minimum size=3pt,
    inner sep=0pt,
]
